\section{Residual tame eigenvalues on an original member}

% BEGIN R86G1
\paragraph{R86G1. The cuspidal point and its inertia polynomial.}
Fix an original member $D=w.D(Z)$ with the layer-point guards, a location
$i$ of coefficient characteristic $\ell$, and a point parameter satisfying
the hypotheses of $\mathrm{g3p}$. Cuspidality makes its list
$\Sigma=[(r,\pi,1)]$. Its point supplies a unitary representative $P$ whose
base-change list contains $(r,\pi\otimes\eta^{-1},1)$. Here the cuspidal
embedding, its twist and localizations, and the characters are the ones
on $D$ itself, by B85U1 and the member's character frame.
In particular the automorphic representation $\Pi$ of $\mathrm{GL}_r$
represented by $\pi$ is cuspidal over the CM field $E$.

The base-change test makes $\Pi\otimes\eta^{-1}$ conjugate self-dual.
The restriction of $\eta$ to real-field ideles is the quadratic
character of $E/E^+$, so $\eta(x\bar x)=1$ by the norm law. Thus
$\eta^c=\eta^{-1}$ and untwisting gives $\Pi^c\simeq\Pi^\vee$.
At each embedding, the exponent $s$ of $\eta^{-1}$ is $1/2$ or
$-1/2$ by the infinity condition on the same frame. The shifted
infinitesimal multiset of $\Pi$, obtained by adding $s$ to each entry,
is a submultiset of the weight of $P$ by the infinity identity in
its base-change test. This last multiset has distinct entries in
$\mathbb Z+(n_V-1)/2$ by the algebraic coefficient weight.
Since $n_V=r+t$ and $s-t/2\in\mathbb Z$ for odd $t$, the
infinitesimal entries of $\Pi$ are distinct and in
$\mathbb Z+(r-1)/2$. This is regular algebraicity in the
normalization of B85U1.

Let $\rho$ be \emph{any} integral representation in the test.
Consider the characteristic-zero attachment $r_{\ell,\iota_i}(\Pi)$
with parameter normalized by $|\det|^{(1-r)/2}$ as in A72F8.
At a cofinite set of unramified places its characteristic polynomial
equals that of $\rho$ on the source's Frobenius elements.
Indeed $\mathrm{IsRepOfSig}$ reads the polynomial of the normalized
local parameter there. A72F7 and the good-support condition of the
cuspidal embedding identify the weak reading with classical local
Langlands. Both sides evaluate the same Weil element. When spherical
Satake coordinates are given at geometric Frobenius, evaluating at
arithmetic Frobenius inverts the entire list on both sides.
Chebotarev makes these Frobenius conjugacy classes dense. The
coefficients of both polynomials are continuous in
$\overline{\mathbb Q}_\ell$ after applying the specified coefficient
embedding to $\rho$, so they agree on every element of $G_E$.

At $w\nmid\ell$ the semisimplified local-global comparison suffices:
the Weil semisimplification of this attachment is the parameter of
$\Pi_w|\det|^{(1-r)/2}$ up to semisimplification
(I.~Varma, \emph{Local-global compatibility for regular algebraic
cuspidal automorphic representations when $\ell\ne p$},
Forum Math. Sigma 12 (2024), e21, Theorem~1).
On inertia the normalizing twist is trivial. Passing to
semisimplification preserves characteristic polynomials; the
unipotent exponential in the construction of the Weil--Deligne
representation also preserves them there, since it commutes with
the finite inertia matrix. A72F7 and B85U1 thus identify
$\operatorname{charpol}(\rho(u))$, in complex coefficients, with
$D.\mathrm{naut.inePol}(\mathrm{locG}(\pi,w),u)$ for every inertia
element $u$ at a designated type place $w\nmid\ell$.
% END R86G1

% BEGIN R86G2
\paragraph{R86G2. Simultaneous reduction at the designated generator.}
Put $g=g_k\ne\ell$, $w_+=\mathrm{plGw}(k,v)$,
$\gamma=\mathrm{tg}(k,v)=\mathrm{gam0}(i,w_+)$ and let $L$ be the
polynomial label of $\Sigma$ there. Write $S$ for its multiset of
roots in $\overline{\mathbb F}_g$. By full degree and the original
inertia-polynomial reading, R86G1 gives
$\prod_{a\in S}(X-\tau(a))$ for the polynomial on $\gamma$,
where $\tau$ is the member's multiplicative Teichm\"uller lift.
Every $a$ is nonzero since the matrix is invertible and
$\tau(0)=0$ (multiplicativity and injectivity on units force this).
Put $j=j(i)-1$. The member's separate test on its first line
$\chi_1$ gives $\chi_1(\gamma)=\tau(c)^{\delta}$ in the same
embedding, for $c=\chi_T(k,v,j)\in\mathbb F_g^\times$ and
$\delta\in\{1,-1\}$.

To reduce both formulas, embed $O_i$ in $\mathbb C$ as prescribed
and take its integral closure $B$ there. The roots of unity $\tau(a)$
belong to $B$. Choose a prime of $B$ over the maximal ideal of
$O_i$ by lying over. Its residue field is algebraic over the residue
field of $O_i$, so its embedding into $\overline{\mathbb F}_\ell$
can be chosen to extend the specified embedding $\mathrm{kFbar}_i$.
Composing $\tau$ with this reduction gives one multiplicative map
$J:\overline{\mathbb F}_g\to\overline{\mathbb F}_\ell$.
For $m$ prime to $\ell$, two distinct $m$th roots of unity have
different reductions: the factorization of $X^m-1$ in $B[X]$ would
otherwise have a repeated root after reduction, contrary to its
nonzero derivative at every root. Thus $J$ is injective on the
nonzero elements of order prime to $\ell$. Every $\ell$-power root
reduces to one. For a nonzero $a$ let $a_{\ell'}$ denote its
prime-to-$\ell$ component, the unique factor of order prime to
$\ell$ whose quotient from $a$ has $\ell$-power order. This is
precisely the CRT exponent in $\mathrm{primeToPart}$. Coefficientwise
reduction now gives
\[
 \operatorname{charpol}(\bar\rho(\gamma))
       =\prod_{a\in S}(X-J(a_{\ell'})),\qquad
 \bar\chi_1(\gamma)=J(c_{\ell'})^\delta.
\]
Here $\bar\rho=\mathrm{redF}(D,i,\rho)$ and $\bar\chi_1$ uses
the prescribed residue embedding. All multiplicities survive reduction.

By M84T9 the full degree of $L$ makes the local parameter trivial
on wild inertia, giving polynomial $(X-1)^r$ there in R86G1.
Each wild image of $\bar\rho$ is therefore unipotent, of
$\ell$-power order. The residual representation is continuous with
open kernel, and its restriction to the pro-$g$ wild group has
finite $g$-group image. Hence this image is trivial.
An arithmetic Frobenius lift $\phi$ acts by the $g$th power on tame
inertia, because the residue field at $w_+$ is $\mathbb F_g$.
Consequently
$\bar\rho(\phi\gamma\phi^{-1})=\bar\rho(\gamma)^g$ as matrices.
For the rank-one character the twist data express its value on each
wild element as a product of a $(g-1)$st root and an $\ell$-power
root. Its residual value has order prime to $g$, and is also in
the finite image of a pro-$g$ group, hence equals one. Thus the
product representation $(\bar\rho,\bar\chi_1)$ kills wild inertia.
% END R86G2

% BEGIN R86G3
\paragraph{R86G3. The same reduction for the two tests at a conjugate place.}
Every $w\mid g$ is either a designated $w_+$ or its CM conjugate,
as $g$ splits completely. At $w_+$ use R86G2 with signs
$(1,\delta)$. At its conjugate $w$ write $s=\mathrm{gam0}(i,w)$
and assume the required tame-generator condition there.
CM conjugation and an inner conjugation adjusting the prolongation
of the valuation transport $s$ to a tame generator $s'$ at
$w_+$. This uses transitivity on prolongations and the fact that
the field transport preserves inertia and its wild subgroup.

For clarity, the characteristic polynomial of $\rho(s)$ equals that
of $\rho((s')^{-1})$, and $\chi_1(s)=\chi_1((s')^{-1})$.
Indeed at cofinally many split good primes, conjugate
self-duality in R86G1 inverts the unitary Satake roots.
The common normalization multiplies their values at arithmetic
Frobenius by $q^{(1-r)/2}$. Thus conjugating $\rho$ has the
polynomials of its dual twisted by $\varepsilon_\ell^{1-r}$,
where $\varepsilon_\ell$ is cyclotomic. For $\chi_1$ the
Frobenius formula has power $(t-1)/2+(1-n_V)/2$; the two
conjugate values of each of $\mu_C$ and $\eta_T$ are mutual
inverses on the unramified split test. Therefore the product
of the two conjugate character values is
$\varepsilon_\ell^{t-n_V}$. These equalities extend by continuity
and Chebotarev: for each finite Galois quotient over $E$ take its
Galois closure over $E^+$ and prescribe a lift fixing $E$. A prime
with that Frobenius splits in $E$, and a choice of place realizes
the desired conjugacy class. At inertia away from $\ell$ the
cyclotomic character is one. Inner conjugation preserves
polynomials and character values, giving the asserted identities
also after reduction.

Write $A=(\bar\rho(\gamma),\bar\chi_1(\gamma))$, $M=\mathrm{ord}(A)$.
Both factors have open kernel and R86G2 kills wild inertia in their
product. The resulting finite image of inertia is cyclic.
Since $s'$ and $\gamma$ both generate its tame quotient
topologically before mapping, their images generate the same cyclic
group. The product value at $(s')^{-1}$ is $A^b$ for one integer
$b$ prime to $M$, including the trivial group of order one.
Let $\Theta$ be a multiplicative automorphism of
$\overline{\mathbb F}_g^\times$ acting by the $b$th power on
$M$-torsion. Explicitly decompose each torsion element into its
primary components. On its $p$-primary component raise to $b$ if
$p\mid M$, and act identically otherwise. These powers are
invertible on the specified components, so this gives the required
order-preserving automorphism. Extend it by $\Theta(0)=0$.

Each $J(a_{\ell'})$ is an eigenvalue of the first component of
$A$, and $J(c_{\ell'})^\delta$ is its second component's value.
Their orders divide $M$. Injectivity of $J$ on prime-to-$\ell$
torsion implies that the orders of $a_{\ell'}$ and $c_{\ell'}$
also divide $M$. Thus replacing $J$ by $J\circ\Theta$ raises
\emph{all} these values to the same $b$th power and preserves the
required injectivity. Triangularizing the matrix over the residue
algebraic closure shows that taking the $b$th power raises every
eigenvalue to $b$ with its multiplicity, regardless of Jordan
blocks. Hence the displayed formulas of R86G2 hold at $s$ with this
one new map and with the same $v,j$ and signs $(1,\delta)$.
Together with the Frobenius matrix equality of R86G2, this proves
both clauses of $\mathrm{ResidualTameEigenvalueCondition}$ for
every $i,\Sigma,\rho,k$ in the original test. It supplies
$\mathrm{GaloisPreconditions.g3p}$ on each selected member.
% END R86G3

% BEGIN R86G4
\paragraph{R86G4. Full degree on every point parameter and cuspidality on a good one.}
At a designated $w_+$ the local parameter of any $P\in\mathrm{AV}(i)$
is tame even when $g=\ell_i$. To see the latter case, apply the
fixed-vector argument of D85G1 directly. Its level input
$\mathrm{mt.uni}$ ranges over every coefficient place, every row and
each designated $k,v$; the tuple character is tame there also when
$g=\ell_i$, since both the reference character and the multiplier
are tame. A single-place principal-congruence element has trivial
level character. The nonzero type space of $P$ therefore yields a
nonzero locally fixed vector by the two projections of D85G1. The
depth law applies to its local representation, whose support is good
by the compact-averaging and parabolic-induction argument there.
This is entirely a statement on the complex parameter.

For any $\mathrm{HasParamX}$ witness $P$ the split identity expresses
this tame local parameter as the direct sum of the local strings of
$\Sigma$ and the distinguished string, with the common rank-one
twist $\eta^{-1}$. The member makes $\eta$ unramified at every
finite place. A direct summand of a representation on which wild
inertia acts trivially also has trivial wild action. Untwisting
therefore shows that the entire parameter of $\Sigma$ is tame.
Its label degree is the dimension of its wild invariants by M84T9,
and hence is its full dimension, for all $i,k,v$.
In particular this proves $\mathrm{GlobalDatumDepthZeroCondition}$
on the member for every list at a good point.

Suppose now that $\Sigma$ itself is good. Write its parts as
$(a_h,\pi_h,d_h)$. Positivity in the augmented base-change list
gives $a_h,d_h\ge1$ and $\sum_h a_hd_h=n_V-t=r$.
At each designated place, its label is the product of the
constituent labels to the powers $d_h$: local strings add direct
summands with unramified shifts and the inertia eigenvalues in
M84T9 multiply as polynomials under direct sum.
Every summand is tame, so each factor has degree $a_h>0$.
Squarefreeness forces $d_h=1$. If there were at least two
parts, the smallest $a_h$ would satisfy $1\le a_h\le r/2$,
and its factor would divide the label at \emph{every} designated
place above $g_0$. Goodness supplies a place where that divisor
degree is missing. Thus the list is $[(r,\pi,1)]$.
% END R86G4

% BEGIN R86G5
\paragraph{R86G5. The Hecke polynomial after removing the distinguished string.}
Choose one good point parameter as in R86G4, put
$K=\operatorname{Frac}O_i$, and take the attachment $R$ of $\Pi$
from R86G1 over a finite extension of $K$ in
$\overline{\mathbb Q}_\ell$. We construct the scalar twist explicitly.
The Hecke character $\alpha=\eta^{-1}|\cdot|^{-t/2}$ is algebraic:
the infinity exponents $\pm1/2-t/2$ are integers, since $t$ is odd.
Its $\ell$-adic avatar $\xi$ is unramified at every $w\nmid\ell$
because $\eta$ is unramified and the norm is trivial on finite
units. The construction is the one of M84T10: transport
$\alpha_f^{-1}$ to $\overline{\mathbb Q}_\ell$ and divide by the
embedding powers on the $\ell$-components to make the character
trivial on principal ideles; apply global reciprocity in the
arithmetic-Frobenius convention. For $q=\mathrm Nw$,
\[
       \iota_i\xi(\mathrm{Frob}_w)=\eta_w(\varpi_w)q^{-t/2}.
\]
These values belong to the image of $K$. Indeed the member's
column carries $\mathrm{ValLL0}$ on this frame. Its $\eta$ clause
places $\eta_w(\varpi_w)q^{-1/2}$ in the embedded small field,
and $q^{-(t-1)/2}$ is rational. The member's coefficient embedding
agrees there. Density of Frobenius and closedness of the finite
extension $K$ inside $\overline{\mathbb Q}_\ell$ give
$\xi(G_E)\subset K^\times$. Compactness makes its image under
the valuation a finite subgroup of $\mathbb Z$, hence zero.
Thus the original $O_i^\times$ contains all values of $\xi$.

Set $R'=R\otimes\xi$ and choose over a finite coefficient extension
$K'/K$ a stable lattice. Such a lattice is the span over $O_{K'}$
of all translates of a basis lattice: compactness bounds them
between two scalar multiples of that lattice. On every split
place $w$ in the $\mathrm{SplitOut}$ test the member's $S_i,Q_i$
exclusions imply that $w$ lies above neither a type prime nor the
selected auxiliary prime, and $w\nmid\ell$. The local parameter of
$P$ is unramified by its native unramifiedness clause.
Its base-change identity therefore makes the constituent
$\Pi_w\eta_w^{-1}$ unramified, including zero monodromy.
The classical comparison for the polarized cusp in A72F8 gives
unramifiedness of $R$ there and the normalized Satake polynomial.
The point equality $\mathrm{pt\_loc}$, evaluated at arithmetic
Frobenius $F$, gives
\[
  H_{x,w}(X)=\operatorname{charpol}(R'(F))
                  \prod_{c=1}^{t}(X-\chi_c(F))
\]
in the common coefficient field. In fact the $\Pi$ block in the
augmented list acquires precisely the above scalar, since
$(1-n_V)/2=(1-r)/2-t/2$. The first distinguished root is
$\chi_1(F)$ by the character normalization, and the $k$th
successive shift multiplies it by $q^{-k}$, exactly
$\chi_{k+1}(F)$ for $0\le k<t$. This verifies every factor.
Monic division on $O_i[X]$ shows that the polynomial of $R'(F)$
already has coefficients in $O_i$. Split Frobenius density as in
R86G3 and closedness of $O_i$ give the same assertion for all
$g\in G_E$.

For every residual representation $r_b$ in $\mathrm{IsRhoBar}$,
its defining polynomial times the same distinguished product is
the reduction of $H_{x,w}$. Cancel that monic product and use the
same density, now with locally constant polynomials, to obtain
\[
 \operatorname{charpol}(r_b(g))=
       \operatorname{charpol}(\overline{R'}(g))\quad(g\in G_E)
\]
in $\overline{\mathbb F}_\ell$, using a residue embedding for
$K'$ extending the one for $O_i$.
% END R86G5

% BEGIN R86G6
\paragraph{R86G6. Missing divisor degrees imply the original residual irreducibility test.}
Let $U$ be either the reduction of $R'$ on its stable lattice
or any $r_b$ as in R86G5, extended to $\overline{\mathbb F}_\ell$.
The reduction of $R'$ at inertia away from $\ell$ has exactly the
polynomials of $R$ there since $\xi$ is unramified. Apply the
proofs of R86G1--3 to this stable lattice over $O_{K'}$, using the
extended residue embedding. The designated polynomial and the
two conjugate-place formulas therefore hold also for $U$ by
R86G5. For the designated Frobenius \emph{matrix} identity on an
arbitrary $r_b$, note that its wild inertia polynomial is
$(X-1)^r$ by the same equality. Its kernel is open by assumption.
Its wild image is consequently trivial by the coprime-orders
argument of R86G2, which gives the matrix identity on $r_b$ itself.

The two bad sets $\mathrm{BB}(r,g_k)$ are disjoint in the
original type data. Choose $k$ with $\ell$ outside its bad set,
so $\ell\ne g_k$. At each $v$ above this $g_k$, all eigenvalues
of $U(\gamma_v)$ are distinct: squarefreeness gives $r$ distinct
roots $a$, the good tuple separates their $a_{\ell'}$, and $J$ is
injective on those. If $U$ had a nonzero proper invariant subspace,
take that space or its quotient of dimension $d\le r/2$.
Its $\gamma_v$-eigenvalues form a $d$-element subset of the
full list. The Frobenius matrix identity makes this subset stable
under $g_k$th powers. The bijection $a\mapsto J(a_{\ell'})$ on
the label roots commutes with $g_k$th powers; projection to
prime-to-$\ell$ torsion and $J$ are both multiplicative there.
Consequently the preimage gives a Frobenius-stable $d$-element
subset of the label roots, hence a degree-$d$ divisor over
$\mathbb F_{g_k}$. This holds for every $v$, contrary to goodness
for that degree. This proves absolute irreducibility of $U$.

For \emph{each} $k$ with $g_k\ne\ell$, regardless of its bad set,
the original $S_i$ definition puts every $w\mid g_k$ in $S_i$,
so $\mathrm{gam0}(i,w)$ satisfies the tame-generator test.
The character and matrix polynomials of R86G2--3 use one map
injective on prime-to-$\ell$ torsion. If the first character
were an eigenvalue, the prime-to-$\ell$ part of a label root
would equal that of $\chi_T(k,v,j)$ or its inverse. The two
goodness tests exclude both at every such $\ell$.
Together these results prove $\mathrm{GaloisPreconditions.rhoBar}$:
$\mathrm{HG}$ supplies the good point and list used above, and
the argument covers every residual test representation over its
original residue field. Root exclusion descends because that
field embeds in $\overline{\mathbb F}_\ell$.
% END R86G6

% BEGIN R86G7
\paragraph{R86G7. Integral descent at a good point.}
For each good point $x$ and list $\Sigma$ in the integral residual representation condition,
R86G4--6 apply to this very list. We descend $R'$ from $K'$ to
the specified $O=O_i$. In a basis of the stable $O_{K'}$-lattice,
write $B$ for the $O$-span of its group matrices. This is a
unital algebra because the product of two group matrices is a
group matrix. Put $k=O/\mathfrak m$, $k'=O_{K'}/\mathfrak m'$.
Absolute irreducibility in R86G6 and Burnside's matrix algebra
theorem over $\overline{\mathbb F}_\ell$ give $r^2$ group
matrices $b_1,\ldots,b_{r^2}$ whose reductions form a basis of
$M_r(k')$. In particular they form a basis over $K'$. The trace
pairing on the full matrix algebra is perfect in every
characteristic (dual matrix units are transposes). Hence
$\det(\operatorname{tr}(b_j b_h))$ is a unit in $O_{K'}$.
All entries lie in $O$ by R86G5 applied to products, so this
determinant is a unit in $O$. The coordinates of every group
matrix in the $b_j$ basis are obtained by inverting this Gram
matrix against its traces of products with the $b_h$. They
belong to $O$, proving $B=\bigoplus_j O b_j$.

Reduction and extension to $k'$ now identify
$(B/\mathfrak m B)\otimes_k k'$ with $M_r(k')$. Thus
$B/\mathfrak m B$ is central simple over $k$ (simplicity and
the center can both be tested after this field extension).
The residue field $k$ is finite, so its central simple algebra
is split by the finite division-ring theorem. Lift a rank-one
idempotent of $M_r(k)\simeq B/\mathfrak m B$ to an idempotent
$e\in B$. For example, start from any lift $a$ and iterate
$a\mapsto 3a^2-2a^3$. This fixes $a$ modulo $a^2-a$ and
squares the ideal of the error. The sequence converges in $B$,
which is finite free over the complete DVR $O$. Its limit works
also in residue characteristic two.

The left ideal $Be$ is a direct summand of a free module over
$O$; reduction identifies it with a column module of dimension
$r$, so it is free of rank $r$. The action map
$B\to\operatorname{End}_O(Be)$ is an isomorphism by reduction
and Nakayama (both free ranks are $r^2$). After extending to
$K'$ its simple module of dimension $r$ is isomorphic to the
original matrix representation, since a full matrix algebra
has one simple module up to isomorphism. We have thus obtained
$R'_O:G_E\to\mathrm{GL}_r(O)$ with the same characteristic
polynomials. Its continuity follows from the Gram-matrix
coordinate formula and continuity of the original matrices:
the coordinates lie in the closed subfield $K$ with its induced
topology. Its generic fiber is absolutely irreducible by the
full matrix algebra just computed.

Set $\rho=R'_O\otimes\xi^{-1}$ with $\xi$ from R86G5.
This gives $\mathrm{IsRepOfSig}$ on the original lattice, since
over $K'$ it is the ordinary attachment $R$. The twist $\xi$
is trivial on inertia at every type prime away from $\ell$,
and $\rho\otimes\xi=R'_O$ satisfies the Hecke polynomial
identity at every $\mathrm{SplitOut}$ place by R86G5.
These are exactly the three clauses of the original
$\mathrm{GlobalDatumResidualRepCondition}$, for the specified
$x$ and $\Sigma$. Thus the $\mathrm{rho}$ and $\mathrm{dz}$
inputs of the row conditions also follow on original members.
% END R86G7
